\documentclass[12pt]{article}
\usepackage[T1]{fontenc}
\usepackage[utf8]{inputenc}
\usepackage{lmodern}
\usepackage{textcomp}
\usepackage{enumitem}
\usepackage{geometry}
\geometry{margin=1in}
\begin{document}
\begin{center}
Answer Key - Philosophy - Undergraduate Level Assessment 3 \
\small (Answers are ordered exactly as in the question paper.)
\end{center}

\section*{Multiple Choice}
\begin{enumerate}[label=\arabic*.]
\item \textbf{Answer:} E
\\ \textit{Options:} A) actions that affect only themselves.; B) actions performed under duress.; C) actions that do not violate anyone's rights.; D) actions that benefit others.; E) what is due to factors beyond their control.
\\ \textit{Note:} Nagel argues that individuals cannot be held morally accountable for actions or choices that result from circumstances or factors outside their control, as these external influences significantly undermine the notion of personal responsibility.
\item \textbf{Answer:} A
\\ \textit{Options:} A) argument from authority; B) ad hominem attack; C) begging the question; D) red herring; E) significance
\\ \textit{Note:} The argument from authority fallacy occurs when someone asserts that a claim is true based solely on the authority or credibility of the person making the claim, rather than providing substantive evidence or reasoning to support it.
\item \textbf{Answer:} C
\\ \textit{Options:} A) art and beauty.; B) love and hatred.; C) virtue and vice.; D) success and failure.; E) truth and falsehood.
\\ \textit{Note:} Hume argues that our emotional responses to images of happiness or misery are crucial for moral judgments, suggesting that a person who lacks such emotional engagement would also be indifferent to concepts of virtue and vice, which rely on these feelings for their significance.
\item \textbf{Answer:} A
\\ \textit{Options:} A) Thomas Aquinas; B) Catherine of Siena; C) Teresa of Avila; D) Hildegard of Bingen; E) John of the Cross
\\ \textit{Note:} Thomas Aquinas argued that evil is not a substance but a privation of good, which allows for a clearer understanding of divine love by contrasting it with the absence of good.
\item \textbf{Answer:} C
\\ \textit{Options:} A) we have unlimited resources to address environmental issues.; B) human satisfaction is the only measure of environmental success.; C) we can measure in some way the incremental units of human satisfaction.; D) we can precisely predict the environmental impact of our actions.; E) the environment is solely a human construct.
\\ \textit{Note:} The correct answer reflects Baxter's assertion that any environmental goal must be quantifiable in relation to the benefits it provides to human satisfaction, emphasizing the necessity of measuring the outcomes of our actions to assess their environmental impact.
\end{enumerate}

\section*{True / False}
\begin{enumerate}[label=\arabic*.]
\setcounter{enumi}{5}
\item \textbf{Answer:} False
\\ \textit{Options:} A) True; B) False
\\ \textit{Note:} Nathanson argues that all individuals, regardless of their actions, retain their inherent dignity and rights to humane treatment, emphasizing the moral obligation to uphold these rights even for those who have committed serious crimes.
\item \textbf{Answer:} False
\\ \textit{Options:} A) True; B) False
\\ \textit{Note:} Anscombe argues that Sidgwick recognizes the importance of consequences in moral evaluation, suggesting that actions should be judged not only by intentions but also by their outcomes.
\item \textbf{Answer:} True
\\ \textit{Options:} A) True; B) False
\\ \textit{Note:} Rawls' theory of justice emphasizes the importance of the "difference principle," which advocates for social and economic inequalities to be arranged so that they benefit the least advantaged members of society, thereby supporting the idea that students with fewer native assets should receive equivalent attention and resources as their more advantaged peers.
\item \textbf{Answer:} True
\\ \textit{Options:} A) True; B) False
\\ \textit{Note:} Stevenson posits that ethical judgments serve primarily as persuasive tools aimed at shaping the beliefs and actions of others, rather than merely reflecting personal moral views.
\item \textbf{Answer:} False
\\ \textit{Options:} A) True; B) False
\\ \textit{Note:} The statement is false because a valid argument requires not only the correct application of a general rule but also that the premises logically lead to the conclusion, which is not guaranteed in this case.
\end{enumerate}

\section*{Long Form Response}
\begin{enumerate}[label=\arabic*.]
\setcounter{enumi}{10}
\item \textbf{Answer:} In Aquinas's natural law theory, intrinsic goods are fundamental to understanding moral reasoning and human flourishing. These goods are inherently valuable and align with human nature, guiding individuals toward actions that promote well-being and fulfillment. By recognizing and pursuing intrinsic goods-such as knowledge, friendship, and virtue- individuals can make moral choices that reflect their true purpose and contribute to the common good. This framework encourages a holistic approach to ethics, where the pursuit of these goods fosters both personal and communal flourishing, thereby reinforcing the interconnectedness of moral reasoning and human experience. Ultimately, intrinsic goods serve as essential benchmarks in Aquinas's moral philosophy, emphasizing the importance of aligning actions with human nature for a flourishing life.
\\ \textit{Note:} correct because it accurately identifies intrinsic goods as essential components of Aquinas's natural law theory that guide moral reasoning and promote human flourishing by aligning individual actions with their true purpose and the common good.
\item \textbf{Answer:} Velleman argues that the option of euthanasia can harm patients by imposing a significant burden on them to justify their continued existence. This expectation can lead individuals to feel that their lives must meet certain standards of value or productivity, forcing them to rationalize their existence in a way that may not be consistent with their personal experiences or feelings. As a result, patients may experience added psychological distress and a sense of inadequacy if they perceive themselves as a burden to others. Ultimately, this pressure can undermine the intrinsic value of life itself, suggesting that the availability of euthanasia may inadvertently devalue those who are suffering and lead to a harmful societal narrative regarding vulnerability and worth.
\\ \textit{Note:} correct because it captures Velleman's argument that the availability of euthanasia places an undue psychological burden on individuals to justify the worthiness of their lives, which can lead to feelings of inadequacy and increased distress.
\end{enumerate}

\vfill
\noindent\textit{End of Answer Key}
\end{document}